\documentclass[11pt,a4paper]{article}
\usepackage[utf8]{inputenc}
\usepackage[indonesian]{babel}
\usepackage{amsmath,amsfonts,amssymb,amsthm}
\usepackage{geometry}
\geometry{a4paper, margin=2cm}

\title{\textbf{Penyelesaian QUIZ 2 Metode Matematika}}
\author{Dosen: Suhud Wahyudi}
\date{}

\begin{document}

\maketitle

\noindent\textbf{Sifat: Tutup Buku. Waktu: 60 menit}

\vspace{0.5cm}

\begin{enumerate}
  % ---- SOAL 1 ----
  \item Selesaikan Persamaan Beda: $U_{n+2} - 4U_n = 2^n$

        \textbf{Penyelesaian:} \\
        Persamaan ini adalah persamaan beda linear non-homogen orde dua. Solusi umumnya adalah penjumlahan dari solusi homogen ($U_n^{(h)}$) dan solusi partikular ($U_n^{(p)}$), yaitu $U_n = U_n^{(h)} + U_n^{(p)}$.

        \textit{Solusi Homogen:} \\
        Persamaan homogennya adalah $U_{n+2} - 4U_n = 0$. \\
        Persamaan karakteristiknya: $r^2 - 4 = 0$, yang memberikan akar $r_1 = 2$ dan $r_2 = -2$. \\
        Maka, solusi homogennya adalah:
        \begin{equation*}
          U_n^{(h)} = C_1 2^n + C_2 (-2)^n
        \end{equation*}

        \textit{Solusi Partikular:} \\
        Bagian non-homogen adalah $f(n) = 2^n$. Karena basis $2$ merupakan salah satu akar tunggal dari persamaan karakteristik, maka kita tebak solusi partikularnya dalam bentuk:
        \begin{equation*}
          U_n^{(p)} = A \cdot n \cdot 2^n
        \end{equation*}
        Substitusi $U_n^{(p)}$ ke dalam persamaan awal $U_{n+2} - 4U_n = 2^n$:
        \begin{align*}
          A(n+2)2^{n+2} - 4(A \cdot n \cdot 2^n) & = 2^n                                                                \\
          A(n+2)(4 \cdot 2^n) - 4An(2^n)         & = 2^n                                                                \\
          4A(n+2) - 4An                          & = 1 \quad \rightarrow \text{(membagi kedua ruas dengan } 2^n\text{)} \\
          4An + 8A - 4An                         & = 1                                                                  \\
          8A                                     & = 1 \implies A = \frac{1}{8}
        \end{align*}
        Maka solusi partikularnya adalah $U_n^{(p)} = \frac{1}{8} n 2^n = n 2^{n-3}$.

        \textit{Solusi Umum:}
        \begin{equation*}
          U_n = C_1 2^n + C_2 (-2)^n + n 2^{n-3}
        \end{equation*}

        \vspace{0.5cm}

        % ---- SOAL 2 ----
  \item Selesaikan PD: $(1 - x)y'' - 4xy' - 2y = 5x - 3$ disekitar $x = 0$

        \textbf{Penyelesaian:} \\
        Karena koefisien $y''$ bernilai tak nol di $x=0$ ($1-0 \neq 0$), maka $x=0$ adalah titik biasa (ordinary point). Kita gunakan deret pangkat:
        \begin{equation*}
          y = \sum_{n=0}^{\infty} a_n x^n, \quad y' = \sum_{n=1}^{\infty} n a_n x^{n-1}, \quad y'' = \sum_{n=2}^{\infty} n(n-1) a_n x^{n-2}
        \end{equation*}
        Substitusikan deret-deret tersebut ke dalam persamaan diferensial:
        \begin{equation*}
          (1-x)\sum_{n=2}^{\infty} n(n-1) a_n x^{n-2} - 4x \sum_{n=1}^{\infty} n a_n x^{n-1} - 2\sum_{n=0}^{\infty} a_n x^n = 5x - 3
        \end{equation*}
        Jabarkan dan pisahkan summasinya:
        \begin{equation*}
          \sum_{n=2}^{\infty} n(n-1) a_n x^{n-2} - \sum_{n=2}^{\infty} n(n-1) a_n x^{n-1} - \sum_{n=1}^{\infty} 4n a_n x^n - \sum_{n=0}^{\infty} 2a_n x^n = 5x - 3
        \end{equation*}
        Ubah indeks agar semua suku menggunakan pangkat $x^k$:
        \begin{align*}
          \sum_{k=0}^{\infty} (k+2)(k+1) a_{k+2} x^k - \sum_{k=1}^{\infty} (k+1)k a_{k+1} x^k - \sum_{k=1}^{\infty} 4k a_k x^k - \sum_{k=0}^{\infty} 2a_k x^k & = 5x - 3
        \end{align*}
        Kumpulkan koefisien $x^k$ untuk masing-masing $k$:

        \textit{Untuk $k=0$:}
        \begin{align*}
          2(1) a_2 - 2 a_0 & = -3 \implies 2a_2 - 2a_0 = -3 \implies a_2 = a_0 - \frac{3}{2}
        \end{align*}

        \textit{Untuk $k=1$:}
        \begin{align*}
          3(2) a_3 - 2(1) a_2 - 4(1) a_1 - 2 a_1 & = 5 \\
          6a_3 - 2a_2 - 6a_1                     & = 5
        \end{align*}
        Substitusi nilai $a_2$:
        \begin{equation*}
          6a_3 - 2\left(a_0 - \frac{3}{2}\right) - 6a_1 = 5 \implies 6a_3 - 2a_0 + 3 - 6a_1 = 5 \implies a_3 = \frac{1}{3}a_0 + a_1 + \frac{1}{3}
        \end{equation*}

        \textit{Untuk $k \geq 2$:}
        \begin{equation*}
          (k+2)(k+1) a_{k+2} - k(k+1) a_{k+1} - 4k a_k - 2a_k = 0
        \end{equation*}
        \begin{equation*}
          (k+2)(k+1) a_{k+2} - k(k+1) a_{k+1} - (4k+2) a_k = 0
        \end{equation*}
        Sehingga didapatkan relasi rekursi untuk $k \geq 2$:
        \begin{equation*}
          a_{k+2} = \frac{k(k+1) a_{k+1} + 2(2k+1) a_k}{(k+2)(k+1)}
        \end{equation*}

        Jadi, solusi deretnya dapat dituliskan sebagai:
        \begin{equation*}
          y(x) = a_0 + a_1 x + \left(a_0 - \frac{3}{2}\right) x^2 + \left(\frac{1}{3}a_0 + a_1 + \frac{1}{3}\right) x^3 + \dots
        \end{equation*}
        dimana $a_0$ dan $a_1$ adalah konstanta sembarang dari \textit{Initial Value}.

        \newpage

        % ---- SOAL 3 ----
  \item Selesaikan dengan metode Frobenius Persamaan Diferensial berikut: $2x^2y'' + x(2x + 1)y' - y = 0$ disekitar titik $x = 0$

        \textbf{Penyelesaian:} \\
        $x=0$ adalah titik singular reguler. Kita selesaikan dengan metode deret Frobenius:
        \begin{equation*}
          y = \sum_{n=0}^{\infty} c_n x^{n+r}, \quad y' = \sum_{n=0}^{\infty} (n+r) c_n x^{n+r-1}, \quad y'' = \sum_{n=0}^{\infty} (n+r)(n+r-1) c_n x^{n+r-2}
        \end{equation*}
        Substitusikan ke dalam persamaan diferensial awal:
        \begin{equation*}
          2x^2 \sum_{n=0}^{\infty} (n+r)(n+r-1) c_n x^{n+r-2} + x(2x+1) \sum_{n=0}^{\infty} (n+r) c_n x^{n+r-1} - \sum_{n=0}^{\infty} c_n x^{n+r} = 0
        \end{equation*}
        \begin{equation*}
          \sum_{n=0}^{\infty} 2(n+r)(n+r-1) c_n x^{n+r} + \sum_{n=0}^{\infty} 2(n+r) c_n x^{n+r+1} + \sum_{n=0}^{\infty} (n+r) c_n x^{n+r} - \sum_{n=0}^{\infty} c_n x^{n+r} = 0
        \end{equation*}
        Satukan suku dengan pangkat $x^{n+r}$:
        \begin{equation*}
          \sum_{n=0}^{\infty} \left[2(n+r)(n+r-1) + (n+r) - 1\right] c_n x^{n+r} + \sum_{n=0}^{\infty} 2(n+r) c_n x^{n+r+1} = 0
        \end{equation*}
        Perhatikan bahwa $2(n+r)(n+r-1) + (n+r) - 1 = (2n+2r+1)(n+r-1)$.
        Dengan menyamakan indeks $k+1 = n \implies k = n-1$, deret kedua menjadi:
        \begin{equation*}
          \dots + \sum_{n=1}^{\infty} 2(n+r-1) c_{n-1} x^{n+r} = 0
        \end{equation*}

        \textit{Persamaan Indisial (koefisien untuk $n=0$ di deret pertama disamakan dengan nol):}
        \begin{align*}
          [2(0+r)(r-1) + r - 1] c_0 & = 0                                            \\
          (2r+1)(r-1)               & = 0 \implies r_1 = 1, \quad r_2 = -\frac{1}{2}
        \end{align*}

        \textit{Hubungan Rekursi untuk $n \geq 1$:}
        \begin{equation*}
          (2n+2r+1)(n+r-1) c_n + 2(n+r-1) c_{n-1} = 0 \implies c_n = \frac{-2}{2n+2r+1} c_{n-1}
        \end{equation*}

        \textbf{Kasus 1: $r_1 = 1$}
        \begin{equation*}
          c_n = \frac{-2}{2n+3} c_{n-1}
        \end{equation*}
        Maka nilai konstanta-konstantanya (jika $c_0 = a$):
        \begin{equation*}
          c_1 = -\frac{2}{5}a, \quad c_2 = \frac{4}{35}a, \quad c_n = \frac{(-2)^n}{3 \cdot 5 \cdot 7 \dots (2n+1)} a
        \end{equation*}
        Solusi deret pertama $y_1(x) = a x \sum_{n=0}^{\infty} \frac{(-2)^n}{(2n+1)!!} x^n$.

        \textbf{Kasus 2: $r_2 = -1/2$}
        \begin{equation*}
          c_n = \frac{-2}{2n + 2(-1/2) + 1} c_{n-1} = -\frac{1}{n} c_{n-1}
        \end{equation*}
        Maka nilai konstanta-konstantanya (jika $c_0 = b$):
        \begin{equation*}
          c_1 = -b, \quad c_2 = \frac{1}{2}b, \quad c_3 = -\frac{1}{6}b, \quad \dots \quad c_n = \frac{(-1)^n}{n!} b
        \end{equation*}
        Karena $\sum_{n=0}^\infty \frac{(-1)^n}{n!} x^n = e^{-x}$, maka solusi deret kedua adalah $y_2(x) = b x^{-1/2} e^{-x}$.

        \textbf{Solusi Umum:}
        \begin{equation*}
          y = y_1 + y_2 = a x \sum_{n=0}^{\infty} \frac{(-2)^n}{(2n+1)!!} x^n + b \frac{e^{-x}}{\sqrt{x}}
        \end{equation*}

\end{enumerate}

\end{document}